\begin{apartats}

\apartat{2,5}
Calculeu els coeficients $a$, $b$, $c$ i $d$ de la funció $f(x)=ax^3+bx^2+cx+d$ si sabem que l'equació de la recta
tangent a la gràfica de la funció $f$ en el punt d'inflexió $(1,0)$ és $y=-3x+3$ i que la funció té un extrem relatiu
en el punt de la gràfica d'abscissa $x=0$.

\begin{solucio}
Si calculem les funcions derivades successives, tenim $f'(x)=3ax^2+2bx+c$ i $f''(x)=6ax+2b$. Si anem traduint en
termes de $f$ i de les seves derivades les condicions de l'enunciat, obtenim:
\begin{itemize}
\item $f$ passa per $(1,0)\rightarrow f(1)=0\rightarrow a+b+c+d=0$ (equació 1);
\item $(1,0)$ és un punt d'inflexió $\rightarrow f''(1)=0\rightarrow6a+2b=0$ (equació 2);
\item el pendent de la tangent en $x=1$ és igual a $-3\rightarrow f'(1)=-3\rightarrow3a+2b+c=-3$ (equació 3);
\item $f$ té un extrem en $x=0\rightarrow f'(0)=0\rightarrow c=0$ (equació 4).
\end{itemize}
Resolent el sistema format per les equacions (1), (2), (3) i (4), obtenim la solució del problema:
\[
\boxed{a=1,\quad b=-3,\quad c=0,\quad d=2}.
\]

\textit{Pauta oficial:} 0,5 per la traducció de cadascuna de les quatre condicions i 0,5 per la resolució final del
sistema.
\end{solucio}

\end{apartats}
